% 写作提示：把结果放回文献、任务目标和适用范围里讨论，并写明局限。
\section{讨论}
\label{sec:discussion}

\subsection{结果讨论}
方案 C 的综合得分最高，主要因为完整性和表达两项领先；方案 B 在准确性上更好。若课程更重视可以核对的依据，就不该只按\tabref{tab:scores}的最后一列排序，而应提高准确性的权重，并同步修改\eqref{eq:score}和\figref{fig:results}。

\subsection{局限与改进}
示例只有三个方案和三项指标，没有误差、重复测量或外部对照。分数是事先设定的，不能当成实验发现。正式报告应说明样本范围和材料缺口，并估计增加指标或改变权重之后，排序可能如何变化。

\begin{rem}
同一指标在正文、表格和插图中应使用相同的名称和精度。
\end{rem}
